% MT023 / explanation — canonical editable content.
先检查类别比例：1000人中只有50名患者，950名非患者。若模型永远预测阴性，准确率仍是\(950/1000=95\%\)，但50名患者一个也找不到。因此“95\%准确率”不能说明筛查有效。


有两个不同问题。\textbf{类别不平衡}是人数悬殊，训练目标容易被多数类主导；\textbf{错误代价不对称}是漏掉患者比多做一次后续CT更严重。处理人数问题与处理决策偏好不能只写成同一句“多收数据”。


训练时可提高阳性样本权重，或在训练折内重采样，使少数类的错误不会被淹没。划分训练与验证时采用分层抽样，尽量让各部分都有阳性样本；不要先复制少数类再整体划分，否则同一原样本的副本可能跨到训练与验证，制造虚高成绩。


决策时可降低判阳性的阈值，用验证数据权衡漏诊与额外检查。最终测试至少报告混淆矩阵，并计算
\[
\text{召回率/敏感度}=\frac{TP}{TP+FN},\qquad
\text{特异度}=\frac{TN}{TN+FP}.
\]


\(TP,FN\)分别是查出的和漏掉的患者，\(TN,FP\)分别是正确放行和误报的非患者。补充例子：查出45名患者、误报95名非患者，则\(FN=5,TN=855\)，召回率90\%，特异度90\%，准确率也是90\%。虽然低于全判阴性的95\%，它却找到了45名患者，体现了任务真正关心的能力。


阈值与模型都在验证阶段确定，测试集留作最终独立评价。


\textbf{出处：} 2021A期中 Q10，原卷第4页（10分）；2023A期中 Q10，原卷第7页（10分）


\textbf{答案依据：} 2021A期中 Q10 随卷／配套参考答案；条件与解释经复核。
